%% =====================================================================
%% Project Ka-Nova, IEEE Access resubmission
%% Manuscript Access-2026-37966, submission UUID
%% b9e6f9b3-ed48-4565-836d-5633d52b89d3
%%
%% Compile with the official IEEE Access LaTeX template class file
%% (ieeeaccess.cls) from the IEEE Author Center. The class file is NOT
%% vendored in this repository; download the current "IEEE Access LaTeX
%% Template" bundle and place ieeeaccess.cls beside this file.
%%
%%   pdflatex ka_nova_access_2026
%%   bibtex   ka_nova_access_2026
%%   pdflatex ka_nova_access_2026
%%   pdflatex ka_nova_access_2026
%%
%% HOUSE RULES ENFORCED IN THIS SOURCE, do not relax without asking:
%%  1. No em dashes anywhere, in prose or in comments.
%%  2. No emoji anywhere.
%%  3. No standardised effect size (no Cohen's d) for any metric.
%%  4. No number is invented or back-filled. Every quantity not yet
%%     measured appears as \todo{PENDING CAMPAIGN: ...}.
%%  5. Causal language stays scoped to "within the simulation".
%% =====================================================================

\documentclass{ieeeaccess}

\usepackage{cite}
\usepackage{amsmath,amssymb,amsfonts}
\usepackage{algorithmic}
\usepackage{graphicx}
\usepackage{textcomp}
\usepackage{xcolor}
\usepackage{booktabs}
\usepackage{url}

%% ---------------------------------------------------------------------
%% PENDING marker. Every one of these must be resolved or deliberately
%% converted to a stated limitation before the manuscript is uploaded.
%% Set \pendingvisibletrue for the internal review copy and false only
%% when the list is empty.
%% ---------------------------------------------------------------------
\newif\ifpendingvisible
\pendingvisibletrue
\newcommand{\todo}[1]{%
  \ifpendingvisible\textcolor{red}{\textbf{[\MakeUppercase{#1}]}}\fi}

\begin{document}

\history{Resubmission of manuscript Access-2026-37966, originally
submitted 30 July 2026, rejected with invitation to resubmit 31 August
2026.}

%% ---------------------------------------------------------------------
%% TITLE
%%
%% The submitted title was "A Hybrid Generative Agent-Based Model with
%% Runtime Constitutional Validation for Comparative Governance
%% Simulation". It no longer describes the paper: the comparative
%% outcomes claim is withdrawn, and the mechanism audit plus the
%% instrument validation are now the primary content. A title change is
%% visible to returning reviewers and is justified explicitly in the
%% response preamble.
%%
%% DECISION NEEDED FROM K.H.: docs/writing_plan.md section 2 recommends
%% "Runtime Constitutional Validation in a Generative Agent-Based Model:
%% Article-Level Enforcement, Substrate Invariance, and a Mechanism
%% Audit". That recommendation predates the 2026-11-24 instrument audit.
%% "Substrate invariance" is now derived rather than discovered and is
%% demoted in the claim set, so naming it in the title over-promises.
%% The title below substitutes the instrument validation for it.
%% ---------------------------------------------------------------------
\title{Encoding Constitutions as Runtime Constraints in a Generative
Agent-Based Model: Article-Level Enforcement, a Mechanism Audit, and an
Instrument Validation}

\author{
\uppercase{Kaung Htet}\authorrefmark{1},
\IEEEmembership{Graduate Student Member, IEEE},
\uppercase{Samsul Jahith}\authorrefmark{1},
\uppercase{Devyani Patil}\authorrefmark{1},
\uppercase{and Md Saifullah Razali}\authorrefmark{1}}

\address[1]{School of Physics, Engineering and Computer Science,
University of Hertfordshire, Hatfield AL10 9AB, U.K. (delivered at PSB
Academy, Singapore)}

\tfootnote{\todo{PENDING: funding statement, or an explicit declaration
that the work received no external funding.}}

\markboth
{K. Htet \headeretal: Encoding Constitutions as Runtime Constraints in a
Generative Agent-Based Model}
{K. Htet \headeretal: Encoding Constitutions as Runtime Constraints in a
Generative Agent-Based Model}

\corresp{Corresponding author: Kaung Htet (e-mail: kh26aab@herts.ac.uk).}

%% =====================================================================
%% ABSTRACT
%% =====================================================================
\begin{abstract}
Constitutional design is normally evaluated through historical case
study and legal-doctrinal reasoning, methods that cannot execute
counterfactuals. This paper presents Ka-Nova, a hybrid generative
agent-based model in which the substantive provisions of a national
constitution are encoded as a typed, validated parameter set and
enforced at run time by a four-layer Constitutional Validation and
Error-Suppression (CVES) pipeline that intercepts every elite decision
and writes an append-only, article-attributed audit record. Seven
constitutions are encoded: a designed Meritocratic Federal Union (MFU),
the 2008 Constitution of Myanmar, and five historically enacted texts
(South Korea, Thailand, Pakistan, Chile and the German Basic Law). The
paper makes three claims. First, within a controlled pair sharing
population, shock sequence, external layer and generative substrate, an
identical elite decision stream is differentially constrained by two
article-level encodings: the same act, a coup signal emitted by an
elite agent, is intercepted and annulled with article attribution under
one encoding while no encoded provision renders it invalid under the
other, with both counts recoverable per decision from the shipped log.
Second, we report a mechanism audit of our own model that states, with
file and line, which published quantities are produced by constitutional
content and which by engine-level constants, and we report as a finding
that the engine does not discriminate between structurally similar
guardianship constitutions, together with the reason and with why the
obvious remedy would be circular. Third, we report an instrument
validation in which an adversarial audit of our own pipeline found and
fixed a class of defects that pre-determined the model's output,
including a parser that discarded the language model's decision for four
of the seven encodings. We quantify the sensitivity of our own published
enforcement counts to the counting rule, and we withdraw and replace
quantities that the audit showed were not supported. No standardised
effect size is reported anywhere in this paper: four previously
headline separations are accounted for by modelling constants, and we
name each. The generative layer is positioned as a control condition and
as an interpretability and provenance mechanism rather than as a
determinant of aggregate outcomes.
\todo{PENDING CAMPAIGN: one sentence giving the headline enforcement
counts from the re-run, and one giving the collapse-timing result, both
to be inserted only once the campaign reports.}
\end{abstract}

\begin{keywords}
Agent-based modelling, computational social science, constitutional
simulation, generative agents, large language models, model validation,
multi-agent systems, reproducibility, runtime validation, sensitivity
analysis.
\end{keywords}

\titlepgskip=-15pt
\maketitle

%% =====================================================================
\section{Introduction}
\label{sec:intro}
%% =====================================================================

\PARstart{T}{he} design of a constitution is among the highest
consequence decisions a polity makes, yet the tools for evaluating a
constitution before it is adopted remain overwhelmingly qualitative.
Comparative constitutional scholarship reasons from historical precedent
and doctrinal analysis \cite{acemoglu2012why,lijphart2012patterns,
elkins2009endurance}. It cannot execute a constitution, cannot run the
same country twice under two founding documents, and cannot project the
compounding consequences of a single structural choice such as whether
the military holds reserved legislative seats, or whether an
anti-corruption body reports to an independent court or to a
commander-in-chief.

This paper asks whether a constitution can be treated as executable
code, and then asks the harder question that the first version of this
work did not: once it is executable, how much of what the model produces
is actually attributable to the encoded text, and how would we know.

The first question is answered by construction. If the substantive
provisions of a founding document, merit thresholds for public office,
veto chambers and their supermajority requirements, resource-allocation
splits, the independence of the judiciary, the constitutional status of
a coup, can be encoded as machine-readable parameters, then a
scenario-neutral engine can run them. We present Ka-Nova, a hybrid
generative agent-based model in which an 11,000-agent rule-based citizen
population, a chain-reasoned large-language-model elite layer and a
stochastic external world execute fifty simulated years under a given
encoding, and in which every elite decision passes through a four-layer
Constitutional Validation and Error-Suppression (CVES) pipeline before
it can influence the simulation.

The second question is the substance of this revision. Between the
initial submission and this version we conducted an adversarial audit of
our own instrument. That audit established three things that change what
this paper can claim, and we report all three rather than only the
favourable one.

\begin{enumerate}
\item \textbf{A single mechanism, visible in two lines of the engine,
produces the separation we previously attributed to constitutional
content.} The corruption trajectory of a guardianship scenario is
produced by a chain-type branch in which the military arm reads no
constitutional field at all. We have promoted the relevant coefficients
to per-constitution fields so that the branch is named, disclosable and
reachable by sensitivity analysis, and we state plainly that this does
not make structurally similar constitutions separate, and why sourcing a
per-country drift rate from each country's own observed corruption
series would be circular.
\item \textbf{Our previously published enforcement counts were derived
by regular expression over free-text reasoning, because the structured
decision was never logged.} Three defensible counting rules give three
materially different answers from the same shipped file. We publish the
sensitivity table, the strict reconciliation, and the mechanical cause
of the residual, and we withdraw the claim that interception was
complete.
\item \textbf{A hardcoded literal in the decision parser silently
discarded the language model's answer for four of the seven encodings.}
Any prior statement about what those four constitutions emit is
therefore not evidence, including statements made in our own internal
record. We report the defect, the mechanism by which it escaped every
test, and the regression coverage that now prevents its recurrence.
\end{enumerate}

Reported in isolation, the first two items read as corrections. Reported
together with the method that found them, they are the contribution.
A modelling instrument whose own output is partly pre-determined by its
implementation is a general hazard in generative agent-based social
simulation \cite{ziems2024can,sagallm2025}, and it is a hazard that is
structurally hard to detect, because the symptom, a clean and strongly
separated result, is indistinguishable from success. We therefore
present the audit as a methodological contribution rather than as an
erratum, while being explicit that an audit is not a validation: it
establishes what the model does not warrant, not that the remainder is
correct.

\subsection{Contributions}
In order of emphasis:

\begin{enumerate}
\item \textbf{The CVES pipeline}, a four-layer runtime validation
mechanism (structural, constitutional, consensus and theory,
statistical) that makes free-form generative elite output usable as
scientific data and produces an auditable, article-attributed
enforcement record keyed to the active constitution's own provisions.
\item \textbf{A mechanism audit of our own model} (Section
\ref{sec:mechanism}) reporting, with file and line, which published
quantities are driven by constitutional content, which by chain-type and
feedback-layer constants, and which are pinned at a variable boundary.
This section also contains the negative result that the engine does not
separate structurally similar guardianship encodings.
\item \textbf{An instrument validation} (Section \ref{sec:instrument}),
a defect audit of the generative path with per-defect statements of what
each defect would have corrupted, and the regression suite that now
covers the path in continuous integration without requiring a GPU.
\item \textbf{Seven machine-readable, article-level constitutional
encodings} released with per-field provenance and explicit
\texttt{CONFIDENT} and \texttt{GUESS} labels, together with a documented
account of which encoded provisions are mechanically reachable in the
present engine and which are encoded but inert.
\item \textbf{A differential-constraint result on a controlled pair},
in which the same elite act is intercepted with article attribution
under one encoding and propagates uncontested under the other, both
recoverable per decision from the shipped audit trail.
\end{enumerate}

We frame Ka-Nova strictly as a methodological and architectural
contribution. The constitutional scenarios are input configurations for
demonstrating and stress-testing the framework, not normative claims
about any polity. Section \ref{sec:threats} states this scope
explicitly and in the strongest terms, since one of the seven encodings
was authored by the first author.

Section \ref{sec:related} situates the work; Section \ref{sec:method}
details the methodology; Section \ref{sec:setup} the experimental
design; Section \ref{sec:mechanism} the mechanism audit; Section
\ref{sec:instrument} the instrument validation; Section
\ref{sec:results} the results; Section \ref{sec:discussion} the
discussion, threats to validity and limitations; Section
\ref{sec:repro} reproducibility; Section \ref{sec:aiuse} the
disclosure of generative-AI use; Section \ref{sec:conclusion}
concludes.

%% =====================================================================
\section{Related Work}
\label{sec:related}
%% =====================================================================

\subsection{Agent-Based Modelling of Social and Political Systems}
Agent-based modelling generates macro-level social phenomena from micro
behavioural rules. Schelling's segregation model \cite{schelling1971dynamic}
showed stark aggregate patterns emerging from mild individual
preferences, and Epstein and Axtell's Sugarscape
\cite{epstein1996growing} established the growing artificial societies
programme in which distributions, migration and conflict arise bottom
up. Methodological syntheses \cite{bonabeau2002agent,axelrod1997complexity,
farmer2009economy,macal2010tutorial,gilbert2005simulation,tesfatsion2006agent}
argue for agent-based modelling wherever heterogeneity, adaptation and
out-of-equilibrium dynamics matter, conditions that characterise
governance under stress. Ka-Nova's citizen layer follows this tradition
and is implemented in Mesa \cite{kazil2020mesa}. What agent-based
modelling has historically lacked is a principled representation of the
deliberative behaviour of a small number of powerful actors whose
choices are not well captured by fixed rules; this is the gap the
generative layer fills.

\subsection{Generative Agents and LLM-Driven Social Simulation}
Large language models as agents in social simulation is a recent,
fast-moving area. Park et al. \cite{park2023generative,park2022social}
demonstrated believable memory-conditioned individual and collective
behaviour. A parallel line probes whether such models can stand in for
human respondents and economic actors \cite{argyle2023out,horton2023large,
aher2023using}, with broader assessments for computational social
science \cite{ziems2024can} and surveys of LLM-empowered agent-based
modelling \cite{gao2024large,xi2023rise,wang2024survey}. Studies of LLM
societies report emergent cooperation and collapse dynamics
\cite{piatti2024cooperate,li2023camel,chen2024agentverse}. This
literature converges on two cautions that motivate our design: generative
agents can be sycophantic or evasive under adversarial prompting, and
their outputs require structured validation to be usable as data.
Chain-of-thought and agentic reasoning methods \cite{yao2023react,
wei2022chain,shinn2023reflexion} inform our conditioned elite chain.

A caution this literature states less often, and which this paper
documents concretely in Section \ref{sec:instrument}, is that the
engineering surface between a language model and a simulation, prompt
construction, response parsing, schema validation and fallback
behaviour, is itself a place where a result can be manufactured without
anyone intending it. Our strongest instrument defect was a single
hardcoded list length in a parser.

\subsection{Multi-Agent LLM Orchestration and Validation}
Multi-agent LLM frameworks coordinate role-specialised agents for
complex tasks \cite{hong2024metagpt,wu2023autogen,guo2024large}. Most
relevant to this work, SagaLLM \cite{sagallm2025} introduces validation
and transaction guarantees for multi-agent LLM planning, reflecting a
growing recognition that generative agents need runtime constraint
enforcement. CVES differs in being keyed to a specific legal document:
its constitutional layer validates each decision against the active
constitution's provisions rather than against task-generic correctness
criteria. Section \ref{sec:mechanism} reports the coverage limits of
that keying honestly, including the two layer-2 rules that are not
reachable under the present decision schema.

\subsection{Over-Refusal in Small Aligned Language Models}
A concern specific to running small, safety-aligned models as elite
agents is over-refusal: the model declines benign, in-simulation prompts
that superficially resemble harmful requests. OR-Bench
\cite{cui2025orbench} establishes that Llama-family models exhibit high
over-refusal rates on safe prompts as a side effect of conservative
safety alignment. Mechanistic analyses further show that over-refusal is
task-dependent and geometrically distinct from harmful-refusal, so that
global refusal-direction ablation in alignment tuning does not cleanly
remove it \cite{maskey2026overrefusal}, and that refusal is mediated by
multi-dimensional polyhedral concept cones rather than a single
direction \cite{wollschlager2025geometry}, which helps explain why a
model may fail to distinguish an in-simulation coup decision from a real
policy violation. This frames our own observation, that Llama 3.2 3B
occasionally refuses governance prompts such as military-supremacy
directives, as a documented alignment limitation rather than a
simulation defect, and it motivates the CVES L1 detector that separates
genuine refusal from a parse failure (Section \ref{sec:cves}).

\subsection{Policy, Institutional and Constitutional Analysis}
Policy-oriented agent-based models simulate markets and endogenous
agent policy \cite{furtado2022policyspace2,dawid2018agent,
geanakoplos2012getting} and, increasingly, generative-agent policy
laboratories \cite{piao2025gplab}. Two political-theory frameworks
anchor Ka-Nova's evaluation layer. Acemoglu and Robinson
\cite{acemoglu2012why,acemoglu2006economic} distinguish inclusive from
extractive institutions; we operationalise an inclusiveness score inside
CVES. Lijphart \cite{lijphart2012patterns,lijphart1977democracy}
characterises consensus versus majoritarian democracy; we operationalise
a consensus score. Institutional economics \cite{north1990institutions,
ostrom1990governing} and constitutional-endurance scholarship
\cite{elkins2009endurance}, together with large constitutional corpora
and governance indices \cite{ccp2023constitute,coppedge2026vdem16,
marshall2020polity5,law2013sham}, inform the encoding.

Ka-Nova's departure is procedural: rather than coding constitutions for
statistical correlation with observed outcomes, it executes them forward
under identical conditions. That procedural stance carries an obligation
that Section \ref{sec:mechanism} discharges: if a quantity is produced
by the engine rather than by the encoding, the method's own logic
requires us to say so. Table \ref{tab:positioning} positions Ka-Nova
against representative systems on four capabilities; the inclusion
criteria are given in Section \ref{sec:conclusion}.

\begin{table}[!t]
\caption{Positioning of Ka-Nova against representative simulation
systems.}
\label{tab:positioning}
\centering
\small
\begin{tabular}{@{}lcccc@{}}
\toprule
System / work & Popn. ABM & LLM agents & Const. encoded & Runtime const. valid. \\
\midrule
Sugarscape \cite{epstein1996growing}      & Yes     & No  & No  & No \\
Generative Agents \cite{park2023generative} & Yes   & Yes & No  & No \\
PolicySpace2 \cite{furtado2022policyspace2} & Yes   & No  & No  & No \\
Generative policy lab \cite{piao2025gplab}  & Partial & Yes & No & No \\
GovSim \cite{piatti2024cooperate}           & No    & Yes & No  & No \\
SagaLLM \cite{sagallm2025}                  & No    & Yes & No  & Task-level only \\
Ka-Nova (this work)                         & Yes   & Yes & Yes & Yes (article-level) \\
\bottomrule
\end{tabular}
\end{table}

\subsection{Case Selection and the Myanmar Context}
Myanmar provides a sharp test case. Its 2008 Constitution entrenches
military prerogatives, a guaranteed bloc of unelected legislative seats,
a military-appointed share of key ministries, and provisions permitting
suspension of fundamental rights under a state of emergency
\cite{myanmar2008constitution,crouch2019myanmar}, and its post-2021
trajectory offers a real-world reference for phenomena the model
reproduces, including the self-defeating dynamics of internet
suppression as citizens migrate to circumvention tools
\cite{accessnow2023keepiton}.

The five additional encodings were selected to vary structure rather
than to maximise outcome spread. South Korea 1987 and the German Basic
Law are civilian-supremacy texts with, respectively, a recent transition
from military rule and an explicit militant-democracy design. Thailand
2017, Pakistan under the Eighth Amendment, and Chile under the 1980
text as operative in its elected-Congress period are guardianship
designs in which an unelected or partly unelected organ retains a
constitutional veto over elected politics. Three of these were chosen
precisely because they are structurally similar to one another, which is
the condition Reviewer 2 identified as the discriminating test of the
approach. Section \ref{sec:nearwins} reports the result of that test,
and it is negative.

%% =====================================================================
\section{Methodology}
\label{sec:method}
%% =====================================================================

\subsection{Overview and Design Principle}
Ka-Nova is a discrete-time simulation advancing in one-year steps over a
50-year horizon. Its governing design principle is that the population,
the external world and the dynamical loops are shared across scenarios,
and that a scenario is defined by which constitution object the elite
layer and institutions read at initialisation.

The submitted version of this paper described the engine as
scenario-neutral without qualification. That description was too strong,
and Section \ref{sec:mechanism} replaces it with a precise one. The
engine is parameterised by the active constitution on most paths, and it
additionally branches on a single structural axis, \texttt{chain\_type},
with values \texttt{civilian} and \texttt{military}, which selects
between two rule branches and two elite-agent rosters. Several
coefficients inside those branches were, at submission, bare literals
rather than constitutional fields. They have been promoted to
per-constitution fields at identical values, which makes them visible,
disclosable and reachable by sensitivity analysis, and which is verified
byte-identical. The branch itself remains and is reported rather than
removed.

Each simulated year executes an ordered nine-step sequence: (1) the
external layer evolves and may fire shocks; (2) the social-media channel
resolves information openness; (3) the elite agents deliberate in a
chain and each decision is validated by CVES; (4) validated decisions
are broadcast to the shared environment; (5) all citizen agents step;
(6) twelve feedback loops update political, economic and social state;
(7) chain-type rules apply; (8) institutions run their annual sessions;
and (9) state is aggregated into key performance indicators. The model
is documented following the ODD protocol \cite{grimm2006standard,
grimm2010odd,grimm2020odd} to support replication.

\subsection{Constitutional Encoding and the Seven-Constitution Design}
\label{sec:encoding}
The core methodological claim is that a constitution's substantive
provisions can be represented as typed, validated parameters. Each
constitution is a hierarchy of frozen dataclasses mapping to chapters
and articles, each with a \texttt{validate()} method enforcing internal
consistency, for example that merit weights and the resource split each
sum to one, and, since this revision, that the tax schedule is a
contiguous, strictly ascending set of three-tuples with an unbounded top
band. The last of these was added after we established that
\texttt{validate()} did not inspect tax schedules at all, so a malformed
schedule passed validation and then miscomputed silently at run time.

Seven constitutions are encoded (Table \ref{tab:scenarios}). Each
carries per-field provenance comments naming the instrument, the
citation, the year and a confidence tag. Between 35 and 50 percent of
the implemented fields in each of the five new encodings are
self-labelled \texttt{GUESS}, and we report that ratio rather than
suppress it, because it is evidence for one of this paper's findings
(Section \ref{sec:schemafit}).

\begin{table}[!t]
\caption{The seven encoded constitutions. \texttt{chain\_type} is the
single structural axis on which the engine branches; every other
difference between scenarios is a parameter value read from the active
constitution. Year-zero corruption is the V-Dem \texttt{v2x\_corr}
observation for that polity at its own year zero
\cite{coppedge2026vdem16}. The scenario label names the instrument while
the observation is taken at the year the instrument entered force, and
these differ for two encodings: the Korean value is the 1988
observation, because the 1987 constitution entered force on 25 February
1988, and the Chilean value is the 1990 observation for the
elected-Congress period rather than a Junta-era value.}
\label{tab:scenarios}
\centering
\small
\begin{tabular}{@{}llll@{}}
\toprule
Code & Encoding & \texttt{chain\_type} & Year-zero corr. \\
\midrule
A  & MFU (designed)                      & civilian & 0.878 \\
C  & Myanmar 2008                        & military & 0.878 \\
SK & South Korea 1987                    & civilian & 0.458 \\
TH & Thailand 2017                       & military & 0.683 \\
PK & Pakistan, Eighth Amendment 1985     & military & 0.637 \\
CL & Chile 1980, elected-Congress period & civilian & 0.067 \\
DE & German Basic Law through 1968       & civilian & 0.020 \\
\bottomrule
\end{tabular}
\end{table}

Two encoding decisions are worth stating because they were made against
our own interest. First, Myanmar's year-zero corruption was corrected
from 0.72 to 0.878. The 0.72 figure was present at the repository's
initial commit labelled as a V-Dem baseline; it is not a V-Dem value for
any year, and it is in fact the inverted Transparency International
Corruption Perceptions Index score of 28 of 100
\cite{ti2024cpi}. The real V-Dem \texttt{v2x\_corr} value for Myanmar in
2008 is 0.878 \cite{coppedge2026vdem16}. Both members of the controlled
pair were moved together, which preserves the pair, and the correction
makes the designed constitution's simulated corruption outcome
substantially worse than previously published. Second, Scenario A's and
Scenario C's year-zero coup risk is 0.25, which is four to five times
Myanmar's historical base rate computed on the same convention used for
every other scenario (4 successful coups in 78 years, 0.051, from
\cite{powell2011global}). We did not quietly change it. We disclose it,
publish the historical rate beside it, and put \texttt{coup\_risk} into
the sensitivity analysis so that its contribution can be measured rather
than asserted (Section \ref{sec:gsa}).

The scenarios divide into two deliberately labelled experiment classes.
\emph{Class (i)}, the controlled pair A versus C, shares Myanmar's
population, geography, external layer, shock sequence and year-zero
conditions; constitutional text is the only manipulated variable, and
this is verified, with zero mismatches across all eleven external-layer
output keys over a full 50-year same-seed comparison. \emph{Class
(ii)}, the five historically enacted encodings, is a realism-maximised
comparative study in which each country runs with its own real
administrative geography, its own real population distribution and its
own real neighbours and external shock catalogue. Class (ii) is
therefore \emph{not} a controlled experiment, by design, and no claim in
this paper treats a class (ii) difference as caused by constitutional
text alone.

\subsection{Layer 1: Rule-Based Citizen Agents}
The population comprises 11,000 Mesa agents, distributed across each
scenario's own real administrative units using that country's own
population weights, and assigned to census-proportional ethnic groups
and seven behavioural archetypes (civic champion, pragmatic survivor,
regime loyalist, ambitious meritocrat, disillusioned youth, rural
traditionalist, trauma carrier). Each citizen runs a five-stage cycle
each year, perceive, decide, act, learn, life-course, with
recency-weighted memory, and decides on tax compliance, protest,
emigration, bribery and corruption reporting. Four further tiers
complete the institutional structure: tier 2 officials (ministers,
ethnic leaders, analysis council, congress); tier 3 oversight
(anti-corruption inspectorate, constitutional and arbitration judges);
tier 4 foreign actors; and tier 5 institutional agents (central bank,
development fund, tax system). Citizen identity and scheduling use Mesa
while state updates are NumPy-vectorised \cite{harris2020array}.

The population count is uniform at 11,000 for every scenario rather than
scaled to each country's real population. Scaling at a fixed ratio would
have produced a 7.3-fold difference in agent count between the largest
and smallest polity, which would introduce unequal statistical power
across scenarios for reasons unrelated to constitutional structure. Real
population shares are used only to distribute the 11,000 agents across
each country's own geographic units.

One limitation of this layer reaches a published metric and is stated
here rather than in the limitations section, because it governs how the
tax encoding should be read. The income generator produces a hard
ceiling of 5.46 times the poverty line. Real personal-income exemption
thresholds for three of the encoded countries sit above the reachable
range, Pakistan's at 16.1 times, so those schedules are
\emph{encoded and inert} rather than used. We encoded the real statutory
schedules and disclosed the inertness rather than rescaling the brackets
to fit the engine, which would have reintroduced invented content.

\subsection{Layer 2: Generative LLM Elite Agents}
The elite layer models the small set of actors whose behaviour is
deliberative rather than rule-bound. Seven LLM agents per scenario
reason in a fixed chain, each conditioned on a status report and on
prior advisors' natural-language summaries. Under a civilian-supremacy
chain the chain operates under civilian supremacy; under a
military-guardianship chain it follows a military chain of command with
objectives derived from the active constitution. Each agent emits a
structured decision, a budget weight in $[0,1]$, an ethnic-weighting
vector whose length is that constitution's own ethnic-group count, a
boolean coup signal, and a free-text rationale, inside a delimited
block. Agents are served by a locally hosted Llama 3.2 3B model
\cite{dubey2024llama3,touvron2023llama} via Ollama \cite{ollama2023},
orchestrated with LangChain \cite{chase2022langchain}.

The aggregate contribution of the elite layer to citizen trust is
bounded by
\begin{equation}
b = \min(\overline{w}\cdot 0.15,\ b_{\max}),
\qquad \overline{w} = \mathrm{mean}_i(w_i)
\label{eq:coupling}
\end{equation}
where $w_i$ are per-agent budget weights, and the ethnic-weight vector
is normalised to mean one before consumption, so that only its relative
dispersion survives. We draw attention to \eqref{eq:coupling} in the
methodology rather than leaving it to the results, because it is the
reason the substrate-invariance result in Section \ref{sec:llmrole}
holds, and that is better stated in advance than discovered.

A deterministic rule-based decision path provides an equivalent
fallback, which enables the LLM-versus-rule-based comparison.

\subsubsection{Prompt Symmetry}
At submission, the elite prompts for the five historically enacted
constitutions contained instructions telling the model that a coup
signal had no constitutional basis and should not be set. The two
original scenarios contained no such instruction. This is a circular
construction: it tells the model that a coup is impossible under
constitution $X$, observes no coup, and would report that $X$ prevents
coups. It also disables the measurement, because CVES can only intercept
a signal that is emitted. Those instructions have been removed. Every
factual description of what each constitution says is retained,
including every real article reference and the plain statement that the
text contains no coup mechanism; what was removed is every instruction
directed at what the agent should output. Four example decision blocks
whose rationale strings pre-argued the answer were neutralised. After
the change, no scenario carries an output-fixing directive in any of its
seven role prompts, and every military role carries the same latitude
the designed constitution's military role always had. Whether a given
constitution's elites then decline to signal is a measurement, not an
instruction.

\subsection{Layer 3: External World and Information Environment}
The external layer evolves a per-scenario set of real geopolitical
influence vectors, each with a calibrated annual drift, and fires
stochastic shock types. For the controlled pair, the same seeded shock
sequence hits both scenarios, removing external stochasticity as a
confound. For class (ii) each country runs against its own real
neighbours and its own shock catalogue, which is a deliberate shift
toward realism and away from engine symmetry, and which is one reason
class (ii) is not presented as a controlled comparison.

The social-media channel models information openness as an outcome of
elite choice: when the regime attempts suppression, a VPN-floor mechanic
raises a lower bound on effective openness, so effective openness is
\begin{equation}
O_{\mathrm{eff}} = \max(O_{\mathrm{regime}},\ \mathrm{VPN}_{\mathrm{floor}}),
\qquad \mathrm{VPN}_{\mathrm{floor}} \uparrow \text{ per shutdown}
\label{eq:vpn}
\end{equation}
making suppression self-defeating over time, calibrated to Myanmar's
post-2021 experience \cite{accessnow2023keepiton}.

\subsection{CVES: The Constitutional Validation and Error-Suppression Pipeline}
\label{sec:cves}
Every elite decision passes through the four-layer pipeline before
influencing the simulation; all violations and layer scores are logged
to append-only JSONL for audit.

\emph{L1, structural integrity.} Raw output is parsed and validated
against a typed schema: budget weight clamped to $[0,1]$; an
ethnic-weight vector whose required length is the active constitution's
own \texttt{ETHNIC\_GROUPS} count, with element values clamped to
$[0.3,3.0]$; a boolean coup signal; a rationale. On failure the agent is
re-prompted up to three times, after which a role default is substituted
and the failure flagged, so a parse failure is never miscounted as a
behavioural refusal. Section \ref{sec:instrument} reports a defect in
which this layer's own schema was correct but the upstream parser was
not.

\emph{L2, constitutional constraint.} The decision is checked against
the active constitution. The layer currently implements three
article-level rules, and we state their reachability precisely rather
than their count. (a) A coup signal emitted on a civilian-supremacy
chain is a violation; this rule is live for every civilian-chain
encoding, keyed on \texttt{chain\_type} and on the active constitution's
own coup fields, and it is the rule that produces the enforcement record
in this paper. (b) A rights violation under a constitution whose rights
are non-suspendable is a violation; across the seven encodings only the
designed constitution and the German Basic Law encode non-suspendable
rights, and the rule is additionally unreachable at present because no
prompt, schema or parser produces a \texttt{rights\_violation} field.
(c) A budget weight outside $[0,1]$ is a violation; this rule is
structurally redundant, because the parser clips the value upstream, and
we label it as such rather than engineer it to fire. The article labels
attached to these rules are the designed constitution's own numbering.
We therefore do not claim per-article attribution over historically
enacted constitutions; we claim article-attributed enforcement for the
encoding whose article numbering the label belongs to, and we report the
labelling limitation as an open item in Section \ref{sec:limitations}.

\emph{L3, consensus and theory alignment.} Across the seven decisions,
budget weights are clamped toward the ensemble (structured negotiation),
and two theory scores are logged, an inclusiveness score
\cite{acemoglu2012why} and a consensus score \cite{lijphart2012patterns}:
\begin{equation}
I = \frac{\mathrm{mean}(w_{\mathrm{minority}})}{\mathrm{mean}(w_{\mathrm{majority}})},
\qquad
K = 1 - \mathrm{var}(w_{\mathrm{ethnic}})
\label{eq:theory}
\end{equation}
These components are skipped when the active constitution encodes fewer
than two ethnic groups, which is the case for four of the seven
encodings. That is a coverage limit of L3 and is stated as one.

\emph{L4, statistical plausibility.} The combined budget impact is
$z$-tested against a rolling baseline; an outlier is flagged when
\begin{equation}
\frac{|x - \mu_{\mathrm{hist}}|}{\sigma_{\mathrm{hist}}} > 3
\label{eq:l4}
\end{equation}
The pipeline runs in flag mode, logging and retaining the decision, for
every result reported in this paper and in the submitted version. A
strip mode exists and has never been used in any reported result. This
is specified in the project's integration specification, which states
that outliers beyond three standard deviations are to be flagged and
not rejected.

A separate detector distinguishes a true refusal by the language model
from a valid terse decision and from a parse failure.

\subsection{Outcome Metrics}
\label{sec:metrics}
Primary indicators are the corruption index, public-trust index, coup
probability, anti-corruption body effectiveness, ethnic harmony,
stability index and Gini coefficient, each computed annually. Scenario
differentiation is assessed with the Mann-Whitney $U$ test
\cite{mann1947test}, with Shapiro-Wilk \cite{shapiro1965analysis} used
to check normality.

\textbf{No standardised effect size is reported anywhere in this paper.}
The submitted version reported four, and Section \ref{sec:effectsize}
explains why each is withdrawn. Where a separation between scenarios is
complete, we report it as complete and name the constant that makes it
so, rather than converting a constant into a statistic.

A composite north-star progress score is also computed as a weighted sum
of seven terms. Following the audit, the composite is reported
\emph{decomposed term by term for every scenario} and never as a single
headline number, so that a reader can see which terms are
constitution-driven and which are structural. The weights were inherited
from an earlier phase of the project and were never chosen for a
seven-constitution comparison; we neither re-weight them, because any
weights we chose would reintroduce the self-design problem, nor retire
them, because the composite is one of the metrics a reviewer asked us to
resolve rather than remove.

Global sensitivity analysis follows established practice for agent-based
models \cite{tenbroeke2016sensitivity,lorscheid2012opening} and
systematic experimental design, implemented per Saltelli et al.
\cite{saltelli2008global} via SALib \cite{herman2017salib}. Its scope is
stated in Section \ref{sec:gsa} and it is narrower than the submitted
version implied.

%% =====================================================================
\section{Experimental Setup}
\label{sec:setup}
%% =====================================================================

The experiment has two classes, as defined in Section
\ref{sec:encoding}. Table \ref{tab:config} summarises the configuration.

\begin{table}[!t]
\caption{Experimental configuration.}
\label{tab:config}
\centering
\small
\begin{tabular}{@{}ll@{}}
\toprule
Parameter & Value \\
\midrule
Class (i), controlled pair      & A (MFU) vs C (Myanmar 2008) \\
Class (ii), enacted encodings   & SK, TH, PK, CL, DE \\
Runs per scenario, rule-based   & 100 \\
Runs per scenario, LLM          & \todo{pending campaign: final LLM roster and $N$ per scenario} \\
Horizon                         & 50 simulated years \\
Citizen agents                  & 11,000 (all scenarios) \\
LLM elite agents                & 7 per scenario \\
Model / serving                 & Llama 3.2 3B / Ollama \\
Hardware                        & \todo{pending campaign: GPU model and driver of the campaign pod} \\
Integrity controls              & \texttt{use\_llm} verified in run log; per-call \\
                                & fallback counter; connectivity probe; \\
                                & consecutive-fallback circuit breaker \\
\bottomrule
\end{tabular}
\end{table}

The verification and validation protocol follows established guidance
for simulation models \cite{sargent2013verification} and for the
empirical validation of agent-based models \cite{windrum2007empirical}.
Before the full run, a pre-flight suite verifies syntactic integrity,
model initialisation under every encoding, the direction of indicator
differentiation over the horizon, the elite-to-citizen feedback path,
and the external and social-media layers. An adversarial debiasing suite
(blind-scenario, condition-injection and suppression-audit modules) and
a red-team check on suppressed emergence guard against artificially tame
output.

Two integrity controls are new in this revision and are reported because
their absence would have invalidated a paid campaign. First, language
model initialisation now performs a real round-trip at start-up and
aborts the run on failure, because constructing the client opens no
connection and the realistic failure mode, a model that was never
pulled, previously survived initialisation and surfaced only as a 100
percent per-call fallback rate under an LLM filename. Second, a circuit
breaker aborts after 21 consecutive fallbacks, which is three full years
of a seven-agent chain, so that a mid-run server death cannot produce
hours of rule-based rows inside an LLM dataset. Per-call fallbacks are
stamped on the logged row, so the rule-based contamination rate of any
campaign is a single key lookup.

%% =====================================================================
\section{Mechanism Audit}
\label{sec:mechanism}
%% =====================================================================

This section reports what the engine does and does not discriminate. It
requires no campaign data: every statement in it is verifiable against
the published source at the commit cited in Section \ref{sec:repro}, and
most were produced by reading our own code adversarially rather than by
running it.

\subsection{What the Engine Branches On}
\label{sec:branch}
The engine branches on exactly one structural axis,
\texttt{chain\_type}, resolved from a scenario registry at
initialisation. Every other inter-scenario difference is a parameter
value read from the active constitution. This is a weaker statement than
the submitted manuscript's, and the difference matters: a branch is not
a parameter, and a sensitivity analysis over parameters cannot see it.

The most consequential instance is the state-level corruption update in
\texttt{model\_phase3.py}, method \texttt{\_update\_states}. Under a
civilian chain, state corruption is reduced in proportion to
anti-corruption effectiveness; under a military chain, it accumulates at
a fixed drift:
\begin{equation}
c_{t+1} =
\begin{cases}
\mathrm{clip}\big(c_t - \iota_t \cdot \kappa,\ 0,\ 1\big) & \text{civilian chain}\\[2pt]
\mathrm{clip}\big(c_t + \delta,\ 0,\ 1\big) & \text{military chain}
\end{cases}
\label{eq:corruption}
\end{equation}
where $\iota_t$ is anti-corruption effectiveness, $\kappa$ is the
reduction coefficient and $\delta$ is the unchecked drift. At submission
both $\kappa=0.02$ and $\delta=0.01$ were bare literals inside the
branch, and \textbf{the military arm read no constitutional field at
all}. The consequence, which we measured rather than inferred, is that
the three guardianship scenarios produce corruption series that are
parallel straight lines of slope exactly 0.0100 per year, differing only
in their year-zero intercept.

All five coefficients in that method, $\kappa$, $\delta$, and the three
coefficients of the growth update, are now per-constitution fields read
via a lookup with the historical literal as fallback. The promotion was
verified byte-identical across all seven scenarios at matched seeds,
21 of 21 per-run outputs hash-identical. We stress that \textbf{this
does not make the near-twins separate}, for the reason given in Section
\ref{sec:nearwins}. What it does is make the branch named, inspectable,
disclosable and reachable by sensitivity analysis, which is a
precondition for honest reporting and not a result.

A second instance is the anti-corruption effectiveness floor for
civilian chains, which is still a bare literal of the form
$\min(0.30 + \text{year bonus},\ 0.40)$ and which appears at two
separate sites in the engine. We report it here as an open item rather
than promote it silently.

\subsection{The Ordering Defect in the Effectiveness Floor}
That floor was also, at submission, applied in the wrong order. It was
enforced at the end of the state-update method, after the per-state loop
had already consumed the effectiveness value. The institutional
sub-model overwrites that value earlier in the year with its real
computed figure, which is zero before any investigation resolves. Every
constitution except the designed one encodes its own effectiveness
floor, which props the value back up; the designed constitution encodes
none. The result, measured at a fixed seed before the fix, is that
Scenario A read an effectiveness of exactly zero for its first five
years and applied no corruption reduction at all, so its corruption
index was frozen at 0.8976 for precisely those five years. The published
effectiveness column and the number the corruption mechanism consumed
were different quantities.

This is a real behavioural fix and is not byte-identical for Scenario A,
which is why it is committed separately from the coefficient promotion.
It moves the designed constitution's outcome in the favourable
direction, and we therefore flag it explicitly as a change that benefits
the self-designed scenario, so that a reader can weigh it accordingly.
\todo{PENDING CAMPAIGN: the recomputed Scenario A corruption trajectory
under the corrected ordering at the publication $N$. The provisional
single-seed figure measured during the fix is a 24 percent reduction in
year-50 corruption, which is not a publication quantity.}

\subsection{Separation by Construction: The Four Governing Constants}
\label{sec:effectsize}
The submitted manuscript reported four standardised effect sizes for the
A-versus-C comparison and described them as very large. Reviewer 1 asked
whether such separations are mechanically induced by thresholds,
parameter bounds or deterministic relationships embedded in the model.
We investigated all four, and the answer is that each is accounted for
by a modelling constant. We therefore report no standardised effect size
anywhere in this paper.

\begin{table*}[!t]
\caption{The four previously published separations and the constant that
accounts for each. ``S1'' and ``ST'' are first-order and total-order
Sobol indices from the sensitivity analysis described in Section
\ref{sec:gsa}. No standardised effect size is reported in this paper.}
\label{tab:fourconstants}
\centering
\small
\begin{tabular}{@{}lll@{}}
\toprule
Metric & Governing constant & Evidence \\
\midrule
Public trust
  & \texttt{rights.FEAR\_COMPLIANCE\_TRUST\_FLOOR} $=0.10$
  & Sobol ST $=1.00$ on Scenario C; the reported value is the clamp, and \\
  & & its variance across runs is exactly zero. \\
Coup probability
  & civilian-chain coup-signal damping, $0.05$ per signal
  & Sobol indices undefined: the output has zero variance across all 960 \\
  & & samples of the full parameter space. Constant by construction. \\
Stability
  & \texttt{rights.STABILITY\_RIGHTS\_BONUS} $=0.20$
  & Sobol S1 $=0.9908$; the other twelve parameters all sit at ST $\le 0.008$. \\
Corruption
  & the \texttt{chain\_type} branch of \eqref{eq:corruption}
  & The military arm reads no constitutional field; the three guardianship \\
  & & series are parallel lines of slope exactly $0.0100$/yr. \\
\bottomrule
\end{tabular}
\end{table*}

Table \ref{tab:fourconstants} gives the account. Three of the four were
established by global sensitivity analysis; the fourth, corruption, was
established by reading the source, and it could not have been
established by sensitivity analysis, because \texttt{chain\_type} is a
branch rather than a sampled parameter. The corruption case also
contradicts the mechanism narrative given in the submitted manuscript's
discussion, which attributed accumulating corruption under the 2008
encoding to regime-protective oversight capping anti-corruption
effectiveness. The military branch does not read anti-corruption
effectiveness. We withdraw that narrative.

Two consequences follow for how differences between scenarios are
reported in this paper.

First, where a separation is complete and produced by a constant, we
report it as complete, name the constant, and treat it as a check that
the model behaves as specified rather than as a discovery. A clamp that
holds trust at exactly its floor with zero variance across runs is
evidence that the clamp is implemented; it is not evidence about
constitutions.

Second, horizon-matched comparison is mandatory. The submitted
manuscript compared corruption at the end of the horizon for a scenario
that completes the horizon against a scenario that terminates in its
second decade. We report the horizon-matched comparison at the year
nearest the terminating scenario's mean final year, and we additionally
report that the terminating scenario reaches the variable's upper clip
before its own mean termination year, so that the comparison is
compressed from above by a bound rather than by a mechanism.
\todo{PENDING CAMPAIGN: horizon-matched corruption comparison at the
campaign $N$, with the clip year stated for each scenario.}

\subsection{Boundary-Pinned Metrics}
Corruption is pinned at a boundary for three of the seven encodings.
Germany reaches exactly zero, Chile reaches exactly zero, and the 2008
Myanmar encoding reaches the upper clip of one, in each case before the
end of the period over which it is observed. Any figure plotting
corruption across all seven scenarios therefore contains flat lines at
both bounds, and those flat segments carry no information about the
encodings. Every corruption figure in this paper marks the year at which
each series reaches a bound.
\todo{PENDING CAMPAIGN: the bound-arrival year for each of the seven
series at the campaign $N$.}

\subsection{The North-Star Composite, Decomposed}
The composite progress score is a weighted sum of seven terms with
weights summing to one: corruption 0.20, trust 0.15, coup 0.15,
anti-corruption effectiveness 0.15, growth 0.15, ethnic harmony 0.10 and
inequality 0.10. Measured at full publication scale, \textbf{three of
those seven terms land at an identical maximal value for all four
civilian-chain scenarios regardless of which real constitution is
active}. Those three are the coup term, the effectiveness term and the
growth term, which together carry 45 percent of total composite weight.
Two of the three are civilian-chain hardcodes; the third is a ceiling
effect in which the composite's own growth clip meets uncapped
compounding growth. Nearly half of the composite success score is
therefore decided before any constitution is read.

This is the single most important reason the composite is reported
decomposed and never as a headline number, and it is also the clearest
available illustration of what this paper means by a mechanism audit.
\todo{PENDING CAMPAIGN: the seven-term decomposition table for all seven
scenarios at the campaign $N$.}

\subsection{Global Sensitivity Analysis and Its Scope}
\label{sec:gsa}
Four Sobol analyses have been completed, each at 960 model runs,
$n_{\mathrm{sample}}=64$, over a 13-parameter space at full citizen
scale on the rule-based path. Their results are the evidence column of
Table \ref{tab:fourconstants}. Three cautions must accompany any
quotation of them.

\begin{enumerate}
\item \textbf{The analysis is scoped to the designed constitution.}
The sampling bounds were authored around that constitution's own
parameter values. Checked programmatically against every scenario's live
values, those bounds fail to bracket the real configuration for a
majority of parameters on several of the other encodings, and a spot
check found that 61 percent of non-designed-constitution defaults fall
outside the sampled region. A cross-scenario sensitivity table under one
shared bound set would therefore sample regions that do not contain the
scenarios' actual configurations. The engine now enforces this scope at
run time rather than relying on a documentation note.
\item \textbf{Precision at this sample size is marginal.} The confidence
interval on the dominant stability index is wide relative to the
estimate. The qualitative signal, twelve of thirteen parameters at total
order at or below 0.008, is nonetheless unambiguous. Negative
first-order estimates are estimator noise and are reported as not
distinguishable from zero, never as negative sensitivity.
\item \textbf{A zero-variance output is a finding, not a failure.}
Scenario A's coup probability has zero variance across the whole
parameter space, which makes Sobol indices undefined. We report the
metric as constant by construction.
\end{enumerate}

\todo{PENDING: the two remaining sensitivity runs, a restricted-parameter
run on the Myanmar encoding for trust, and a \texttt{coup\_risk} run,
both of which are free to execute and are required before the
\texttt{coup\_risk} disclosure in Section \ref{sec:encoding} can be
quantified.}

\subsection{Schema Fit, and Why the Comparator Magnitudes Cannot Drive Outcomes}
\label{sec:schemafit}
The parameter schema into which all seven constitutions are encoded was
derived from the designed constitution. That is a measurement problem,
and it is measurable. The designed constitution's schema is 1.6 to 1.8
times larger than any real constitution's encoding, and between 35 and
50 percent of each new country's implemented fields are self-labelled
\texttt{GUESS}.

The anti-corruption institution makes the problem concrete. It is the
designed constitution's own invention. The three real analogues are not
the same kind of institution: Thailand's National Anti-Corruption
Commission investigates and prosecutes after the fact; Chile's
Contraloría exercises an ex ante veto, preventing the act before it
occurs; Pakistan under the encoded instrument has no such body at all.
They differ in mechanism, not in magnitude. Collapsing them into a
single effectiveness scalar discards exactly what distinguishes the
constitutions, and the scalar then has to be filled from somewhere.
Inspecting our own provenance comments, we find that two of the three
comparator magnitudes were calibrated against \emph{outcome} measures,
a control-of-corruption governance index and a corruption perceptions
score, and only one was derived from constitutional text.

The consequence is a constraint on the method, not a defect in the
implementation: a magnitude calibrated from observed corruption cannot
be used to drive a simulated corruption trajectory without circularity.

\subsection{The Structurally Similar Constitutions Do Not Separate}
\label{sec:nearwins}
Reviewer 2 asked how well the approach would work for more similar
constitutions. We built that test, with three independently sourced
guardianship encodings, and we report the result, which is negative.

Inside the engine, the three guardianship encodings differ in their
year-zero intercepts and in their encoded effectiveness floors, and
their corruption trajectories are parallel lines of identical slope,
because the drift coefficient $\delta$ in \eqref{eq:corruption} has no
constitutional analogue in any of the three texts and defaults to the
same value for all three. Promoting it to a constitutional field, which
we did, makes this visible; it does not make the encodings separate.

There is an obvious remedy and we rejected it. One could source a
per-country drift rate from each country's own observed corruption
series. That would make the encodings separate, and it would do so by
fitting the model to the outcome the model then claims to reproduce. A
non-circular version would derive drift from encoded constitutional
features, which is new modelling research rather than a revision fix.

One genuine positive counter-example survives and we report it as the
one place the encoding work demonstrably functions: on the military
side, the anti-corruption effectiveness term tracks Thailand's real
encoded floor at roughly nine to ten times the value produced for
Myanmar and Pakistan. The same mechanism behaves correctly when it is
not sitting beside a competing chain-type hardcode.

The honest summary is that the approach as implemented discriminates
between structurally opposite constitutions and does not discriminate
between structurally similar ones, that we can say precisely which line
of code is responsible, and that we can say what it would take to change
that. We report this as a property of the modelling approach rather than
as an implementation shortcoming, because the deeper cause is the
schema-fit problem of Section \ref{sec:schemafit}: real constitutions
are partly characterised by the absence of the designed constitution's
inventions, and where an analogue exists it differs in mechanism rather
than magnitude.

%% =====================================================================
\section{Instrument Validation}
\label{sec:instrument}
%% =====================================================================

A simulation that embeds a language model has an engineering surface
between the two that is not usually treated as a scientific object. This
section treats it as one. We report the defects our own audit found in
that surface, how each was found, and what each would have corrupted had
it not been found. We state up front that this is an audit and not a
validation: it establishes what the instrument did not warrant, not that
the remainder is correct.

\subsection{Why These Defects Survived Every Test}
One structural fact explains the cluster. Development and all automated
testing run on hardware without a local model server; the language model
path executes only on a rented GPU pod. Every gate in the test suite
therefore exercised the deterministic rule-based path, and the
rule-based path never constructs a prompt and never parses a model
response. Three of the defects below live in code that, before this
revision, \textbf{had never executed once outside a paid run}.

Compounding this, the continuous integration pipeline contained no test
steps at all. It built and published the container image that the
compute pods pull at boot, with no check of any kind, so code that did
not compile would still ship. Builds had passed consistently for exactly
that reason: there was nothing they could fail.

\subsection{The Defects}
Table \ref{tab:defects} lists them. Each row names what the defect was,
how it was found, and what it would have corrupted.

\begin{table*}[!t]
\caption{Defects found by the instrument audit, in the order they were
fixed. Commit identifiers refer to the public repository
\cite{kanova2026repo}. ``Would have corrupted'' states the specific
reviewer-facing quantity at risk, not a generic severity.}
\label{tab:defects}
\centering
\small
\begin{tabular}{@{}p{0.6cm}p{4.2cm}p{4.6cm}p{6.0cm}@{}}
\toprule
ID & Defect & How found & Would have corrupted \\
\midrule
D1 & The run identifier was never written. No writer existed anywhere
     for the field, so every logged row recorded run zero.
   & Attempting to partition the published audit trail by Monte Carlo run.
   & The 100 runs per scenario are not separable in the shipped logs, so
     no per-run statistic can be recomputed from them. Directly implicated
     in the reconciliation failure of Section \ref{sec:table6}. \\
D2 & The structured decision was never logged. The logger received the
     full decision object and wrote only its prose rationale.
   & Attempting to count emitted coup signals from a structured field.
   & Every enforcement count in the submitted manuscript had to be
     recovered by regular expression over free text. See Section
     \ref{sec:table6}. \\
D3 & Language-model initialisation failure silently disabled the model
     and continued; constructing the client opens no connection, so a
     server that was never started survived initialisation.
   & Testing against a deliberately dead endpoint.
   & A paid campaign would have exited successfully having written a
     100 percent rule-based dataset under a language-model filename. \\
D4 & Three encodings failed 100 percent of model calls. Their system
     prompts contain a literal JSON example, and the prompt was passed as
     a template, so a literal brace was parsed as a required variable.
   & A circuit breaker added the same day aborted after 21 consecutive
     failures instead of degrading silently.
   & Any campaign including those three encodings would have produced no
     model output at all, reported as success. \\
D5 & Fixing D4 by passing prompt text literally broke the other four
     encodings, which use the escaped brace convention: the model began
     seeing and copying a literal doubled brace.
   & Reading the model's raw output on the pod after the D4 fix.
   & Malformed decision blocks in 28 of 49 role prompts, with downstream
     parse failures. \\
D6 & Continuous integration had no test steps.
   & Reading the workflow definition after D4.
   & Every green build in the project's history was uninformative about
     code quality. \\
D7 & Prompt pre-suppression. The five enacted encodings instructed the
     model that a coup signal had no constitutional basis; the two
     original encodings did not.
   & A systematic read of all 49 role prompts prompted by the principal
     author's instruction that every constitution be treated on the same
     terms.
   & The central measurement. Telling a model that an act is impossible
     under an encoding and then reporting that the encoding prevents the
     act is circular, and it disables the validator, which can only
     intercept a signal that is emitted. \\
D8 & The anti-corruption effectiveness floor was applied after its
     consumer rather than before.
   & Tracing the corruption term's input live.
   & The designed constitution's corruption trajectory, frozen for five
     simulated years. See Section \ref{sec:branch}. \\
D9 & \textbf{The decision parser accepted an ethnic-weight vector only
     if its length was exactly eight, a hardcoded value from one country's
     census.}
   & One encoding running at roughly ten times the expected latency with
     every decision flagged as a parse failure, while its raw model output
     was valid.
   & \textbf{A guaranteed parse failure on every decision for four of the
     seven encodings, with the model's answer discarded and a role default
     substituted, including the coup signal. No prior statement about what
     those four encodings emit is evidence.} \\
\bottomrule
\end{tabular}
\end{table*}

\subsection{D9 in Detail, and What It Invalidates}
D9 is the one with consequences for claims rather than for operations,
so we state it precisely. An earlier fix had generalised the CVES L1
\emph{schema} to read each constitution's own ethnic-group count, but
the upstream \emph{parser} retained a literal length check. For any
encoding whose count is not eight, the model's correctly formed answer
was discarded, the role default was retained, and the schema then
correctly rejected the default for having the wrong length. On
persistent L1 failure the pipeline substitutes role defaults, so the
model's coup signal was discarded along with the rest and replaced by
the default value of false. Four of the seven encodings have an ethnic
group count other than eight; only the two Myanmar-based encodings ever
worked.

The consequence is that observations of zero coup signals from those
four encodings, which we had attributed to prompt pre-suppression (D7),
were not a test of that hypothesis. D7 was real and the fix was correct,
but D9 is a sufficient explanation on its own, and the two were never
separated experimentally. We therefore treat all pre-fix measurements on
those encodings as void rather than as evidence, including measurements
recorded in our own internal notes.

\subsection{Regression Coverage}
The gap that allowed D4, D5 and D9 is now covered by a test that
requires no GPU and no network. It renders the real prompt object for
all 49 role prompts across all seven encodings and asserts that the
rendered text is byte-identical to the source, that the model is never
shown a literal doubled brace, that every example decision block parses
as valid JSON, and that the exact decision block each encoding instructs
its model to emit survives \emph{that encoding's own} L1 parse with its
fields intact. Verified in both directions: it passes on the current
engine and fails on 21 of 49 prompts against the pre-fix construction,
so it would have caught D4 before a pod was rented. Writing it
immediately surfaced a further defect, one role in one encoding
declaring an eight-element ethnic vector against its own encoded count
of zero. Continuous integration now runs this test, the compile sweep,
all seven encodings' validation, the systems diagnostic, the suppression
boundary, the composite decomposition and all seven scenarios end to
end, and the image build depends on that job passing.

\subsection{How This Should and Should Not Be Read}
It should be read as follows: an instrument that can pre-determine its
own output is a general hazard in this class of work, the hazard is hard
to see because its symptom resembles success, and we found ours by
auditing our own code adversarially against our own published claims.
Four of the nine defects were found only because a fail-loud control
added earlier the same day converted a silent degradation into a visible
abort, which is itself a transferable design lesson.

It should not be read as a claim that the instrument is now validated.
Nine defects found is evidence about a search, not about a residual. The
specific items we know remain open are listed in Section
\ref{sec:limitations}.

%% =====================================================================
\section{Results}
\label{sec:results}
%% =====================================================================

\todo{SECTION STATUS: the reconciliation subsection below is complete
and final. Every other subsection is a skeleton whose numbers come from
the campaign. The campaign roster is itself pending: the decision to
exclude the designed constitution from the paid runs and the pre-flight
finding that it is currently the only encoding observed to produce an
enforcement record are in tension, and that tension is unresolved
because the observation predates the D9 fix and is therefore void. Do
not populate any number below from the superseded dataset.}

\subsection{Reconciling the Published Enforcement Record}
\label{sec:table6}
The submitted manuscript reported that CVES intercepted 4,176
constitutionally prohibited coup signals under Scenario A, described
that interception as complete with zero uncorrected signals, and
reported 6,006 permitted coup signals under Scenario C. \textbf{The
completeness claim was not supported and we withdraw it.} This
subsection reports what the shipped audit trail actually supports, and
why the original figures were unreliable in a way that is itself a
methods result.

Because of defect D2, the structured decision was never logged, and the
published counts were recovered by regular expression over free-text
reasoning. A regular expression over prose counts whatever it matches,
and three defensible patterns give three materially different answers
from the same file (Table \ref{tab:counting}).

\begin{table}[!t]
\caption{Counting-rule sensitivity on the shipped audit trail of the
original submission. The same file yields three different answers
depending on the pattern. The published Scenario A figure of 4,570 sits
between the loose and strict counts; the published Scenario C figure of
6,009 is within one percent of the loose count, indicating that the
original analysis counted mentions rather than decisions.}
\label{tab:counting}
\centering
\small
\begin{tabular}{@{}lrr@{}}
\toprule
Counting rule & Scenario A & Scenario C \\
\midrule
Mentions ``coup'' anywhere in the reasoning   & 23,383 & 10,378 \\
A positive coup signal appears anywhere       &  4,866 &  6,069 \\
The delimited decision block itself is positive & \textbf{4,192} & \textbf{4,197} \\
\bottomrule
\end{tabular}
\end{table}

Counted strictly, against the decision actually emitted, the
reconciliation is given in Table \ref{tab:reconciliation}: 4,192 emitted
against 4,176 intercepted, a gap of 16 signals or 0.38 percent, and an
interception rate of 99.62 percent. The gap is therefore an order of
magnitude smaller than the 8.6 percent a loose count implies, but it is
not zero, and the completeness claim remains unsupported.

\begin{table}[!t]
\caption{Strict reconciliation of emitted against intercepted coup
signals, Scenario A, shipped audit trail of the original submission.
Role identifiers are normalised across the two logs, which use display
names in one and role identifiers in the other.}
\label{tab:reconciliation}
\centering
\small
\begin{tabular}{@{}lrrr@{}}
\toprule
Role & Emitted & Intercepted & Gap \\
\midrule
Chief Justice          & 3,732 & 3,690 & $+42$ \\
President              &   269 &   265 & $+4$ \\
Senior General         &   133 &   163 & $-30$ \\
Anti-corruption Director & 57 &    57 & $0$ \\
Central Bank Governor  &     1 &     1 & $0$ \\
\midrule
Total                  & \textbf{4,192} & \textbf{4,176} & \textbf{$+16$} \\
\bottomrule
\end{tabular}
\end{table}

The residual has a mechanical cause rather than being noise, and the
evidence for that is the sign pattern: two roles show more emissions
than interceptions while one shows 30 more interceptions than emissions,
which no under-logging story can produce. When CVES L2 intercepts a
violation it logs the violation and re-prompts the agent, and the
decision finally written to the decision log is the post-correction one.
Where the re-prompt replaced the text, the violation is recorded with no
matching emission. That is precisely the signature observed on the one
role whose prompt grants explicit latitude to signal.

Exact reconciliation on this dataset is impossible, and that is now a
stated limitation rather than an open defect: the coup signal was never
logged as a field (D2) and the run identifier was never written (D1).
Both are fixed. On the re-run, reconciliation becomes a join on run,
year and agent rather than a regular expression, and the residual either
closes or becomes individually attributable.
\todo{PENDING CAMPAIGN: the structured-join reconciliation at the
campaign $N$, replacing Tables \ref{tab:counting} and
\ref{tab:reconciliation} as the primary enforcement record while those
tables are retained as the account of the superseded dataset.}

Two corrections follow for figures carried over from the submitted
manuscript. The Scenario C figure of 6,009 permitted signals must not be
reused; it is a loose count, and the strict figure for the same dataset
is 4,197. And the near-equality of the two strict emission totals, 4,192
and 4,197, is itself the point of the controlled comparison: the same
behaviour, at very different rates per decision, is annulled with
article attribution under one encoding while no encoded provision
renders it invalid under the other.

\subsection{Differential Constraint on the Controlled Pair}
\todo{PENDING CAMPAIGN: the headline enforcement result at the campaign
$N$ with structured logging. Required quantities: emitted and
intercepted coup signals per scenario and per role, the interception
rate with its residual, the per-decision emission rate for each arm, and
the count of enforcement records produced by each encoding in the final
roster. State the article label attached to each interception and whose
article numbering it is.}

\subsection{Scenario Trajectories and Termination}
\todo{PENDING CAMPAIGN: trajectories for the seven indicators across
the final roster at the campaign $N$, with the bound-arrival year marked
on every series that reaches a bound.}

The submitted manuscript reported that the 2008 Myanmar encoding
terminates in coup collapse at a mean simulated year of 16.6 with a
standard deviation of 0.96, in 100 percent of runs. That figure is
superseded by the accumulated fixes and must be recomputed; a free
rule-based survey at full citizen scale and reduced $N$, taken before
several of the fixes in Section \ref{sec:instrument} landed, gives a
provisional mean near 15.3, which moves toward rather than away from the
real-world anchor.
\todo{PENDING CAMPAIGN: the mean final year and standard deviation for
every terminating scenario at the campaign $N$. Until then no collapse
year is stated as final anywhere in this paper. Note for the drafting
team: the mean count of internet shutdowns before collapse is a
different quantity from the collapse year and has been conflated with it
repeatedly; it must never be printed as a year.}

We retain the submitted manuscript's own characterisation of the
historical comparison. The 2008 Constitution was promulgated in 2008 and
the elected government it produced was overthrown by coup in 2021,
approximately 12.7 years later. We report the correspondence between
that interval and the simulated termination interval only as a weak
plausibility check on the encoding, not as validation. The simulated and
historical timelines are not commensurable, because the model abstracts
away the specific triggers, actors and contingencies of the actual
transition. Establishing genuine external validity requires backtesting
against a constitution with a documented post-adoption trajectory, and
the next subsection reports our attempt at that and its outcome.
\todo{PENDING DECISION: the Myanmar anchoring convention is unresolved.
Promulgation to coup is 12.68 years; the operative instrument from 2011
to 2021 is 10.01 years. Whichever is chosen must be applied identically
to Pakistan, where the corresponding intervals are 14.61 and 13.92
years.}

\subsection{Historical Backtesting on a Restricted Subset of Outputs}
Reviewer 1 observed that historical backtesting, even for a restricted
subset of model outputs, would strengthen the evidence. We read that as
licensing fewer \emph{metrics}, not fewer \emph{cases}, and we therefore
report the chosen output set for all five enacted encodings including
where it fails.
\todo{PENDING CAMPAIGN: the backtest table. Rows: all five enacted
encodings. Columns: simulated value, real-world observation, source, and
an explicit verdict including failures. The provisional free-run survey
at reduced $N$ indicates that one case matches its real duration
closely, one survives the full horizon as it did historically, one makes
an out-of-sample prediction beyond its real observation window, and one
over-survives its real duration by roughly 60 percent. None of those
four statements may be printed until the campaign confirms them at the
publication $N$.}

No survival prediction is claimed for any polity. The collapse mechanism
is deliberately symmetric across scenarios, with an identical threshold
and an identical success roll, which preserves the experimental control
and simultaneously limits what the mechanism can represent. We report
that as a scope limit of the mechanism rather than as a result.

\subsection{Information Suppression Dynamics}
\todo{PENDING CAMPAIGN: suppression counts and the terminal VPN floor
for every scenario in the final roster, at the campaign $N$.}
The qualitative mechanism is \eqref{eq:vpn}: within the simulation,
repeated suppression raises the floor on effective openness, so
suppression becomes self-defeating over time, which is consistent with
observed circumvention behaviour \cite{accessnow2023keepiton}.

\subsection{Substrate Comparison}
\todo{PENDING CAMPAIGN: the final-year indicator table under the
language-model elite layer against the deterministic rule-based layer,
for every scenario in the final roster, with the Mann-Whitney result for
each indicator. The submitted manuscript reported a maximum difference
below 0.003 across all indicators; that comparison must be re-executed
because every component of the pipeline downstream of the elite layer
has changed.}

\subsection{The L2 Ablation}
\todo{PENDING CAMPAIGN: the enforce-versus-observe ablation of the
constitutional layer on a civilian-chain scenario. Pre-committed
reporting position, decided before the run and recorded here so that it
cannot be chosen after the result is seen: if the ablation shows that
the suppression of coup probability under the designed constitution is
produced by the validator rather than by the constitutional provision,
that is the finding and it will be reported as such, and the mechanism
narrative in the submitted manuscript's discussion is withdrawn
accordingly. The ablation must run on a civilian chain, because on a
military chain the measured difference is structurally zero: that
encoding reaches no live L2 rule.}

%% =====================================================================
\section{Discussion}
\label{sec:discussion}
%% =====================================================================

\subsection{The Role of the LLM Layer: A Control Condition and a Provenance Mechanism}
\label{sec:llmrole}
Both reviewers independently asked what the generative component
contributes, one framing it as a disjunction between materially
contributing to outcomes and primarily serving as a decision-provenance
and runtime-validation mechanism, the other observing that the layer is
only weakly connected to the main results and asking that its role be
demonstrated more clearly. We take the second horn of the disjunction
and we evidence it, while noting that the reviewer posed a question
rather than endorsing an answer.

\subsubsection{We accept the premise, and we derive it rather than report it}
Within the reported dynamics the generative layer does not determine
aggregate outcomes, and we do not claim it does. More than that: the
invariance is entailed by the coupling design and we state the
derivation rather than presenting the invariance as a discovery. The
layer's total influence on the simulation is three quantities. The first
is a single scalar bounded to $[0, 0.15]$ by \eqref{eq:coupling} and
consumed as a multiplicative term on a transfer that is itself a
fraction of income. The second is an ethnic-weight vector explicitly
normalised to mean one before consumption, so its mean is destroyed by
construction and only relative dispersion survives. The third is a
boolean that a civilian chain suppresses, with a second lock at the
aggregation stage in addition to CVES L2. A bounded scalar, a
mean-normalised vector and a boolean suppressed on four of seven
encodings cannot move aggregate indicators far, and the magnitude of the
observed invariance is the arithmetic consequence of that coupling.

Reporting this as a negative result about generative agents in
governance simulation generally would be an over-claim, and the
submitted manuscript came closer to that than it should have. It is a
property of \emph{this model's} coupling. We report it as such, and we
credit the derivation to our own audit.

\subsubsection{What the layer does contribute}
Three things, in descending order of defensibility.

First, \textbf{it is the only substrate that generates the phenomenon
CVES governs}. A deterministic rule-based elite cannot emit a coup
signal under a civilian-supremacy encoding, because the rule path does
not produce one; the generative elite does, which is what the validator
detects, attributes to an article and annuls, producing an enforcement
record that fixed rules cannot generate. The enforcement record is the
object of study, and it has no rule-based counterpart. This is the
sense, and the only sense, in which the generative layer is necessary to
this paper's central claim.

Second, \textbf{it supplies interpretable decision provenance}. Each
decision carries a natural-language rationale of roughly 260 reasoning
tokens, and the shipped audit trail records, per decision, the
conditions at decision time, the structured decision itself since the
fix for D2, the validation outcome at each of the four layers, and, on
interception, the article cited and the re-prompt that followed. The
layer is a runtime-validation harness with an audit trail.

Third, \textbf{it is a control condition}. Because the rule-based path
produces materially identical aggregate outcomes, a difference between
scenarios cannot be attributed to language-model idiosyncrasy. Within
the simulation it must be attributed to the configuration, and Section
\ref{sec:mechanism} then establishes how much of that configuration is
constitutional content and how much is engine constant. A
single-substrate design could not support the first inference at all.

\subsubsection{The honest cost of this position}
Taking the provenance horn invites a sharper follow-up question: if the
generative component does not determine outcomes, why is it a
contribution at all? Our answer is that the enforcement record is the
contribution and the generative layer is its necessary substrate, and
that a validator with nothing to validate is not a validator. We also
state the limitation plainly: the elite-to-citizen coupling is weak by
construction, so the layer's influence on indicators is small by design,
and strengthening that pathway is the priority for the next iteration of
the engine. That is a design choice we are disclosing, not a result we
are defending.

Finally, D7 in Section \ref{sec:instrument} is directly relevant here
and we place it rather than bury it. Before this revision, the prompts
for five of the seven encodings instructed the model not to emit the
very signal the validator exists to intercept. Any earlier statement
about what the generative layer contributes under those five encodings
was partly a statement about our prompts. That is now fixed, and what
those encodings emit under symmetric prompts is a measurement reported
in Section \ref{sec:results}.

\subsection{Threats to Validity and Scope of Claims}
\label{sec:threats}

\emph{Self-design bias.} Scenario A is an authored constitution. As the
author of the MFU constitution, an inherent self-design bias exists.
Results should be interpreted as demonstrating constitutional parameter
sensitivity rather than normative superiority of any governance
framework. We make no claim that the MFU is optimal or superior; the
paper's claims are methodological and architectural. The MFU serves as
one pole of a deliberately extreme pair chosen to stress-test whether
the pipeline can discriminate constitutions, and the 2008 encoding is a
faithful, real-document anchor. Two further points reinforce this scope
rather than soften it. First, Section \ref{sec:schemafit} reports that
the parameter schema is itself the designed constitution's own
vocabulary, so a real constitution is partly scored on the absence of
that constitution's inventions, which is a measurement artifact and not
a finding. Second, Section \ref{sec:branch} reports a defect whose
correction moved the designed scenario's outcome in its own favour, and
we flag that direction explicitly. The addition of externally sourced
constitutions to the roster is a test of the instrument's coverage and
must not be read as evidence for the designed constitution's merits.

\emph{Confounders: culture, geopolitics, history, economics.} Encoding a
constitution abstracts away the context in which real constitutions
operate. We do not claim that any real country's trajectory is explained
by its constitution alone, and \textbf{none of the causal language in
this paper should be read as extending beyond the simulation}. Within
the controlled pair, the external layer and the population are held
identical, so these factors are controlled for rather than ignored, and
the constitution is the only manipulated variable by design. For the
five enacted encodings this is explicitly not the case: each runs with
its own real geography and its own external context, which is a
deliberate realism choice that forecloses the controlled reading.

\emph{Article labelling.} The article labels written into the
enforcement record are the designed constitution's numbering. An
interception under another encoding therefore carries a label that
belongs to a different document. Reading those labels from the active
constitution requires machine-readable article-number fields that the
encodings do not yet carry, and is identified as required future work.
Until it lands, we do not claim per-article attribution over
historically enacted constitutions.

\emph{Single-model dependence.} Outcomes are conditioned on a single
3B-parameter model; behaviour under larger or different models is a
separate question, and one with a known risk attached, since the
pipeline's re-prompt rate and statistical band are implicitly calibrated
to this model's output variance.

\subsection{Limitations}
\label{sec:limitations}
Beyond the threats above, the following are open and are stated as
limitations rather than resolved.

\begin{enumerate}
\item \textbf{The approach does not separate structurally similar
constitutions} (Section \ref{sec:nearwins}). The route that would make
them separate is circular and was rejected.
\item \textbf{Two of three constitutional-layer rules are unreachable}
under the present decision schema, one because no component produces the
field it tests, one because an upstream clip makes it structurally
redundant (Section \ref{sec:cves}).
\item \textbf{The theory-alignment components are skipped} for the four
encodings with fewer than two encoded ethnic groups.
\item \textbf{Corruption is pinned at a boundary for three of seven
encodings.}
\item \textbf{The real tax schedules are encoded and inert}, not used,
for at least three encodings, because the income generator's ceiling
sits below their real exemption thresholds. We disclose this rather than
rescale the schedules to fit the engine.
\item \textbf{The ethnicity-to-archetype assignment matrix is still
drawn from one country's census for all seven encodings}, so no
cross-country claim involving ethnic harmony is made anywhere in this
paper.
\item \textbf{A rights-violation drop in trust is encoded, differentiated
across constitutions, and structurally unreachable} on civilian chains,
because the chain rules set the triggering flag to false
unconditionally. Real encoding work with no mechanical effect, reported
rather than removed.
\item \textbf{Several year-zero calibration fields remain anchors rather
than observations.} Where a defensible real observation existed we moved
to it and said so; where none did, the field is labelled an explicit
calibration anchor with its rationale. We specifically did not retrofit
a coincidental numerical match as a citation: a brute-force search over
roughly two thousand candidate transformations of a governance panel
hits any two-decimal target by chance, and hitting it is not provenance.
\item \textbf{The container image does not install from the pinned
dependency file}, so the image's installed versions can drift from the
pinned set. This is a known reproducibility defect, deliberately not
changed immediately before a paid campaign, and scheduled for
remediation before the reproducibility package is archived.
\todo{PENDING: confirm this is closed before submission, or keep it as a
stated limitation.}
\end{enumerate}

%% =====================================================================
\section{Reproducibility}
\label{sec:repro}
%% =====================================================================

Reviewer 1 asked specifically for constitutional encodings, parameter
values, random seeds, LLM configuration, prompts, CVES rules and
statistical procedures. Each is addressed below, including where the
answer is a limitation.

\emph{Code and encodings.} The simulation is implemented in Python with
Mesa \cite{kazil2020mesa}, LangChain \cite{chase2022langchain}, NumPy
\cite{harris2020array}, pandas \cite{mckinney2010data}, SciPy
\cite{virtanen2020scipy} and SALib \cite{herman2017salib}; figures use
Matplotlib \cite{hunter2007matplotlib}. All seven constitutional
encodings are released as readable source modules with per-field
provenance comments and explicit confidence tags
\cite{kanova2026repo}.

\emph{Parameter values and provenance.} A provenance table covers every
year-zero field for all seven encodings, naming the instrument, the
observation, the year and the entity, and flagging every field for which
no retrievable observation exists. Governance index point estimates are
re-estimated at each release of their source, so the reproducibility
package ships the frozen data extract rather than a download
instruction.

\emph{Random seeds.} Each Monte Carlo run is seeded from its run
identifier, and the run identifier is now written into every logged row.
It was not, at the time of the original submission, which is defect D1:
the shipped logs of that submission carry a constant run identifier and
cannot be partitioned by run. \textbf{We correct the submitted
manuscript's statement that all runs are seeded}, which was true of the
simulation and false of the audit trail, and we report both facts. In
addition, determinism under the language-model path is bounded rather
than guaranteed: we deliberately do not enable parallel decoding,
because batched decode changes floating-point reduction order and would
put per-seed reproducibility at risk, and we report the model server's
observed seed behaviour explicitly.
\todo{PENDING: the determinism check result, one run identifier executed
twice with the decision logs diffed ignoring timestamps.}

\emph{LLM configuration.} The model, the serving stack, the context
length, the per-call parameters and the exact model digest are recorded
in the run log of every campaign and are reported in the package.
\todo{PENDING CAMPAIGN: model digest, server version, context length and
measured per-call latency of the campaign pod.}

\emph{Prompts.} All 49 role prompts across the seven encodings are
released verbatim. The prompt-symmetry change described in Section
\ref{sec:method} is published as a diff, so that a reader can see
exactly which instructions were removed and verify that every factual
constitutional description was retained.

\emph{CVES rules.} The four layers, the three constitutional rules, the
schema, the clamping ranges and the re-prompt policy are released as
source, and the reachability of each rule is documented in Section
\ref{sec:cves} rather than left to inspection.

\emph{Statistical procedures.} Tests, sample sizes and the sensitivity
analysis configuration are given in Sections \ref{sec:metrics} and
\ref{sec:gsa}, with the scope restriction on the sensitivity analysis
stated explicitly.

\emph{Audit trails.} CVES writes append-only JSONL logs of violations
and layer scores, and the elite decision log records the structured
decision, the conditions at decision time and the run identifier.
Reconstructing the enforcement record from these files is a join rather
than a text search, and the exact reproduction procedure for the
superseded dataset's reconciliation is published alongside it.

\emph{Continuous integration.} The test suite described in Section
\ref{sec:instrument} runs on every commit and gates the container image
build.

%% =====================================================================
\section{Disclosure of Generative-AI Use}
\label{sec:aiuse}
%% =====================================================================

\todo{PENDING FACTUAL DETERMINATION BY THE AUTHORS. The paragraph below
is a template with the correct structure and scope. Every factual claim
in it must be confirmed by the author who produced each artifact before
submission, and any statement that is not confirmed must be removed
rather than softened. In particular the figure provenance question,
whether any of the schematic figures were produced or substantially
modified with a generative tool, has not yet been answered by anyone and
the reviewer asked about it explicitly.}

Generative AI is used inside the object of study: the elite agent layer
is a locally hosted Llama 3.2 3B model, and its role, configuration and
outputs are described in Sections \ref{sec:method} and
\ref{sec:instrument} and released in full. Separately, and as required
by IEEE Access policy, we disclose the use of generative-AI tools in
preparing this manuscript and its artifacts.
\todo{PENDING: state, per artifact class, whether a generative tool was
used: (a) the schematic figures, (b) the data figures, (c) manuscript
prose, (d) source code. For each, name the tool and the extent, and
state that the authors reviewed and take full responsibility for the
content. Data figures are regenerated directly from model output by
scripts released with the package.}

%% =====================================================================
\section{Conclusion}
\label{sec:conclusion}
%% =====================================================================

Ka-Nova demonstrates that a national constitution can be encoded as a
typed, validated, computable parameter set and executed as a multi-decade
controlled experiment, governed at run time by CVES, a four-layer
validation pipeline that makes free-form generative elite decisions
usable as data and produces an auditable, article-attributed enforcement
record. On a controlled pair sharing population, shocks, external layer
and generative substrate, the same elite act is annulled with article
attribution under one encoding and propagates uncontested under the
other, recoverable per decision from the shipped log.

The second half of the contribution is the audit of that instrument by
its own authors. We report, with file and line, which of our published
quantities are produced by constitutional content and which by
engine-level constants; we report that the approach does not separate
structurally similar guardianship encodings and why the obvious remedy
is circular; we quantify the sensitivity of our own enforcement counts to
the counting rule and withdraw the completeness claim that rested on
them; and we report a defect that silently discarded the language
model's decision for four of seven encodings, together with the
regression coverage that now prevents its recurrence. We report no
standardised effect size anywhere, because all four previously published
separations are accounted for by modelling constants that we name.

The systems in Table \ref{tab:positioning} were selected to span the
design space along four capabilities relevant to Ka-Nova: a
population-scale agent layer; LLM-based deliberative agents; encoding of
a constitution as executable parameters; and runtime validation of agent
decisions against those parameters. Candidate systems were drawn from
the agent-based-modelling, generative-agent and computational-policy
literatures surveyed in Section \ref{sec:related}, and a system was
included when it was a peer-reviewed or archival framework that
instantiated at least one of the four capabilities, so that the
comparison exercises each column.

\subsection{Future Work}
Four items are prerequisites rather than extensions. Reading article
labels from the active constitution, which requires machine-readable
article-number fields in each encoding and which is the condition on any
per-article attribution claim over enacted constitutions. Making a
second constitutional rule reachable by wiring a rights-violation field
through schema, parser and prompts, which would make that rule live for
the two encodings with non-suspendable rights. Strengthening the
elite-to-citizen coupling, which is weak by construction and is the
reason the substrate comparison is uninformative about the generative
layer's value. And deriving per-country institutional drift from encoded
constitutional features rather than from observed outcomes, which is the
only non-circular route by which structurally similar constitutions
could separate, and which is new modelling research rather than an
engineering fix.

%% =====================================================================
\section*{Author Contributions}
%% =====================================================================

K.H. conceived and designed the CVES framework. S.J. implemented the
L1-L4 pipeline and conducted agent validation. D.P. encoded the 2008
Myanmar constitutional parameter set and contributed the systematic
literature review underpinning Section \ref{sec:related}. M.S.R.
supervised the work and reviewed the methodology and the claim scope.

%% =====================================================================
\section*{Acknowledgment}
%% =====================================================================

The authors gratefully acknowledge the University of Hertfordshire and
PSB Academy for supporting this research.
\todo{PENDING: confirm with the supervising author whether the Stage-8
acknowledgment revisions constitute a byline change requiring the IEEE
Request for Byline Change form. Do not alter the Author Contributions
statement above in either direction without that determination.}

\bibliographystyle{IEEEtran}
\bibliography{references}

\EOD

\end{document}
